\cotefile{157-2}
\scanpage{51}
\mspage{49}
% archive sheet number: 50
Je présume que non -- qu'il peut arriver que
la longueur minimum requise \inserted{$l(n)$} pour un \deleted{\unsure{chemin}}
chemin dans $\underline{\mathrm{Hom}}(K_n,K'_{\unsure{\tau|t}(n)})$ reliant $f_n$ et $g_n$,
tende vers l'$\infty$ si $n\to\infty$. L'hyp que les
$K'_n$ sont des cat. de Thomason y changerait
quelque chose ? (deux objets homotopes dans
$\underline{\mathrm{Hom}}(X,K')$, si $K'$ de Thomason, sont ils reliés
par une flèche ?).

Bien sûr, on peut prétendre que les catégories
$\underline{\mathrm{Ch}}_{\Sigma,n}(M;\unsure{a},\unsure{b})$ ont été définies ``canoniquement'', et
leur limite inductive $\underline{\mathrm{Ch}}_\Sigma(M;a,b)$ aussi.
Mais la définition était assez arbitraire et abracadabrante -- je suis convaincu que dix mathématiciens,
se proposant de développer cette même
idée des catégories de chemins, dans le même
but que moi, arriveraient à dix définitions
différentes ! Pourtant, il n'est pas exclu que
l'on finisse par trouver un système transitif
d'iso. dans $\unsure{\mathrm{Hot}\,\mathrm{hot}|\mathrm{Cat}\,\mathrm{hot}}$, entre les catégories
de chemins \inserted{``infinis''} \deleted{$K_\infty$} (ou de longueur non précisée)
$K_\infty$ \inserted{$(=K_\infty(M;a,b))$} que nous \unsure{aurons|avons} construites, dans le but
\scanpage{52}
\mspage{50}
% archive sheet number: 51
principal \inserted{(\unsure{disons})} de représenter le foncteur $H$ \uline{sur les
catégories finies}. Et que les isom. obtenus
soient, bien entendu, \unsure{les bons|des bons} dans $\mathrm{Ind}(\mathrm{Cat}\,\mathrm{hot})$.
J'ai l'impression de ne pas avoir encore
axiomatisé à fond la notion même de
chemin, qui après tout a un \deleted{\unsure{signification}} \inserted{sens}
géométrique assez évident et important par
lui-même. Dire \unsure{peu|p.ex.} ce qu'il faut entendre
par ``\inserted{\uline{une}} théorie des chemins dans une catégorie''.
Une telle théorie \inserted{$\unsure{\mathcal T|\tau}$} associerait donc, à
$M,\Sigma,a,b$, un \deleted{\unsure{cat}} système inductif de
catégories $\mathrm{Ch}_{n\,\Sigma}^{\mathcal T}(M;a,b)$. Et les contraintes
géométriques naturelles à imposer à \deleted{\unsure{des}} des
telles théories, en vue \inserted{aussi} de leur utilisation
pour expliciter
\[
\mathrm{Hom}_{M(X)\Sigma_X^{-1}}(a_X,b_X)\simeq
\pi_0\underline{\mathrm{Hom}}\bigl(X,\underline{\underline{\mathrm{Ch}}}^{\mathcal T}(M;a,b)\bigr)
\]
donnerait un \inserted{syst. transitif d'}isom. \deleted{entre les} non seulmt
entre les ind-objets de $\mathrm{Cat}\,\mathrm{hot}$, mais \inserted{peut-être} aussi
entre les $\mathrm{Ch}_{\infty\,\Sigma}^{\mathcal T}(M;a,b)$, pour des
théories différentes. Si on dispose d'un
tel énoncé, la catégorie $\mathrm{Ch}_{\infty\,\Sigma}(M;a,b)$,
et plus \unsure{exactemt|exactement} ses images dans $\mathrm{Cat}\,\mathrm{hot}$ et $\mathrm{Hot}$,
\scanpage{53}
\mspage{51}
% archive sheet number: 52
prendraient un sens \unsure{canonique}\deleted{\unsure{ment}}.

Il y a d'autre part un autre point de vue,
qui peut-être pourrait ``canonifier'' les
cat. $\underline{\mathrm{Ch}}_\infty$ en tant qu'objets de $\mathrm{Cat}\,\mathrm{hot}$
et \inserted{surtout} \deleted{\unsure{de}} $\mathrm{Hot}$ (on passerait à $\mathrm{Cat}\,\mathrm{hot}$ en
appliquant le foncteur $\unsure{\Psi|\psi}$, enveloppe de Thomason-Kan \ldots). En effet, on est intéressé à
construire, en termes de $M,\Sigma$, un ``dérivateur''
\inserted{(ou un prédérivateur)} $\unsure{\mathbb D|\mathscr D}_{M,\Sigma}$ \unsure{surtout} sur le ``domaine'' des catégories
finies. L'étendre à un domaine plus large
est peut-être important dans un ``deuxième
souffle'', mais peut être considéré comme
du luxe dans le premier mouvement \inserted{de} fondation
d'une théorie des dérivateurs. Ceci dit,
on peut se poser des questions du type suivant.

Soit \deleted{$(\mathrm{Hot})_0$}
% slanted
\margin{\uline{NB} $\mathrm{Hot}_0$ est équivalente à une petite catégorie}
\[
\mathrm{Hot}_0\ \unsure{\colon\subset|\subset}\ \mathrm{Hot}
\]
la sous-catégorie pleine de $\mathrm{Hot}$, définie par les
cat. finies -- appelons ces types d'homotopie ``de type
fini'' (ou de \uline{présentation} finie ?). On a un
foncteur évident
\scanpage{54}
\mspage{52}
% archive sheet number: 53
\[
\begin{aligned}
\mathrm{Hot}&\longrightarrow {\mathrm{Hot}_0}^{\wedge}=\underline{\mathrm{Hom}}\bigl(\mathrm{Hot}_0^{\unsure{\circ|o}},(\mathrm{Ens})\bigr),\\
X&\longmapsto\widetilde X=\bigl(\unsure{\xi|\zeta}\mathrel{\unsure{\mapsto|-}}\mathrm{Hom}_{\mathrm{Hot}}(\xi,X)\bigr)
\end{aligned}
\]
Ce foncteur est \uline{conservatif}, puisque un hom
$X\mathrel{\unsure{-|\to}}X'$ dans $\mathrm{Hot}$ est un iso. s'il induit des
iso. pour les foncteurs $\widetilde X\to\widetilde X'$, et $=$ seulmt
pour leurs valeurs pour les sphères $S^n$ ($n\geq 0$).
(Sauf erreur. \deleted{\unsure{C'est en tous cas}} Y \unsure{a}-t-il des canulars,
parce que le \deleted{\unsure{défi}} critère d'équivalence faible par
les $\pi_i$ ne \deleted{\unsure{marche}} \inserted{s'applique} pas directement, vu l'absence
de points bases. \unsure{Ça|Sa} ne doit pas être très sérieux
\ldots) \margin{?\quad\rule[-2ex]{0.4pt}{6ex}} Mais est-il peut-être $=$ \uline{pleinement fidèle} ?
De sorte que $\mathrm{Hot}$ se réaliserait comme sous-catégorie pleine de $ {\mathrm{Hot}_0}^{\wedge}$, contenant les
foncteurs représentables ? \deleted{\unsure{S'il en était}}
\scanpage{55}
% archive sheet number: 54
\begin{flushright}
\unsure{HOT|HDT}\\
VI
\end{flushright}
\scanpage{56}
% archive sheet number: 55
Définition des $W$-équiv. sur $\mathrm{Cat}/S$\qquad p. 41

\unsure{Condition|Conditions} sur \unsure{locale|local.} fond. $\unsure{W|\mathrm{tal}}$\qquad p. 43-45 et p. 1,2

Th.\inserted{-scholie} récapitulatif \unsure{1|I}\qquad p. 70-81

Th.-scholie 2\qquad p. 108-113\qquad avec $\mathrm{Ho}_W^{\mathrm{lc}}(S),\mathrm{Ho}_{W_0}(S)$

beaux diagrammes :\qquad p. 8,\unsure{75|76},175

Nouvelle version théorie HOT :\qquad p. 171-205

(résumé de la situation)

% printed form
% margin note written sideways
\margin{MidiGraf Montpellier 21999}
\scanpage{57}
\mspage{1}
% archive sheet number: 56
Le dérivateur HOT
\margin{14.11.90\\ (Les premières feuilles, d'il y a deux semaines, ont été depuis, dépassées par \uline{VII})}

\deleted{Il faut définir}
Je définis $\mathrm{Hot}$ à partir de la catégorie de
modèles
\[
M=\mathrm{Cat},\qquad W\subset\mathrm{Fl}(\mathrm{Cat}),
\]
$W$ étant en principe les $\mathrm{Hot}$-équivalences. Mais je
vais dans la mesure du possible axiomatiser plus,
en supposant donné un ensemble $W$ plus
général, appelé ``localisateur fondamental'',
(\unsure{basic localizer|basic localiser}), p.ex. $W=\deleted{\unsure{\circ}}W_{\unsure{\mathbb D|\mathscr D}}$, où $\mathbb D$ est un
dérivateur déjà donné (mais en supposant que
son domaine soit $(\mathrm{Cat})$ tout entier). Nous
allons noter au fur et à mesure de quelles
propriétés de $W$ nous aurons besoin. Aux fins de
référence ultérieure, je vais rappeler ici celles que
j'ai dégagées dans P.S.

\uline{L 1}\quad Saturation de $W$ \deleted{\unsure{$\mathbb D$}}
\[
\left\{\begin{aligned}
&\text{a) les identités dans }W\\
&\text{b) Si deux parmi }gf,f,g\text{ dans }W,\text{ le troisième aussi}
\end{aligned}\right.
\]
\diagram{\begin{tikzcd}[ampersand replacement=\&,column sep=small]
X\arrow[r,"f",dash]\&Y\arrow[r,"g",dash]\&Z
\end{tikzcd}}
c) Si $gf$ et $fg$ sont dans $W$, alors $f$ et $g$ aussi
\diagram{\begin{tikzcd}[ampersand replacement=\&]
X\arrow[r,shift left,"f"]\&Y\arrow[l,shift left,"g"]
\end{tikzcd}}

\diagram{\begin{tikzcd}[ampersand replacement=\&,column sep=small,row sep=small]
\&\text{\uline{L 2}}\arrow[d,Rightarrow]\&
\parbox[t]{.64\linewidth}{$\Delta^1$ est univ. $W$-asph. i.e. $\Delta^1\times X\mathrel{\unsure{-|\to}}X$ est dans $W$ pour tout $X$ dans $\mathrm{Cat}$. Cela implique}\\
\&\text{\uline{L 2 bis}}\&
\parbox[t]{.64\linewidth}{$\Delta^1$ est $W$-\unsure{asph.} (C'était L3$'$ dans P.S.)}\\
\&\text{\uline{L 3}}\&
\parbox[t]{.64\linewidth}{Si $X$ a un objet final, $X$ est $W$-asphérique.}\\
\&\&\parbox[t]{.64\linewidth}{Je note ici variantes renforcées\quad $\big|$ Il y a un L3$'$ dual de L3}\\
\&\text{\uline{L 3} bis}\arrow[uu,Rightarrow]\&
\parbox[t]{.64\linewidth}{Si $X$ est $\mathrm{Hot}$-asphérique, il est $W$-asphérique}\\
\&\text{\uline{L 3} ter}\arrow[u,Rightarrow]\arrow[uuuuu,Rightarrow,bend left=48]\arrow[dl]\&
\parbox[t]{.64\linewidth}{Si $X$ est $\mathrm{Hot}$-asphérique, il est universellement $W$-asphérique, i.e $X\times Y\xrightarrow{\mathrm{pr}_2}Y$ est dans $W$ pour tout $Y$ dans $\mathrm{Cat}$.}\\
\parbox{.13\linewidth}{suffit pour entraîner L2,L3}\&\text{\uline{L 3} quater}\arrow[u,Rightarrow]\&
\parbox[t]{.64\linewidth}{Toute $\mathrm{Hot}$-équivalence est dans $W$}
\end{tikzcd}}
\scanpage{58}
\mspage{2}
% archive sheet number: 57
L 4) Si $f:X\mathrel{\unsure{\to|-}}Y$ est $W$-asphérique (i.e. les $X_{/y}\mathrel{\unsure{-|\to}} Y_{/y}$
dans $W$, i.e. (si on a L 3) si les $X_{/y}$ \unsure{est|sont} $W$-asphérique),
alors $f$ aussi. (Axiome de localisation)
% slanted
\margin{il y a un L 4$'$ dual de L 4}

L 5) (Axiome de fibration) Si
\diagram{\begin{tikzcd}[ampersand replacement=\&,column sep=small,row sep=small]
X\arrow[rr,"f",dash]\arrow[dr]\&\&Y\arrow[dl]\\
\&S\&
\end{tikzcd}}
, \deleted{\unsure{avec}} $X$ et $Y$
fibrées sur $S$ et $f$ cartésien, si pour tout $s\in S$
$f_s:X_s\mathrel{-}Y_s$ est dans $W$, $f$ aussi.
% slanted
\margin{Si $W$ un localisateur, alors les catégories contractiles sont $W$-asphériques, les \unsure{homotopismes|homotopies} (\unsure{au s.} $\unsure{\Rightarrow|=}$) $\in W$, les applications homotopes sont égales dans $(\mathrm{Cat})W^{-1}\overset{\mathrm{déf}}{=}\mathrm{Hot}_W$.}

J'ai supposé tjrs \fbox{L.1, L3, L4} (impliquent
L2) -- c'étaient ces conditions qui définissaient la
notion de ``localisateur'' (fondamental) dans Cat.
D'autre part, \fbox{L.1, L\unsure{2|3}bis, L5} impliquent
L1,2,3,4,5.

\uline{On va supposer ici aussi qu'on ait au moins L1,L3,L4} \inserted{(i.e. $W$ \inserted{un} localisateur)}.
En appliquant la construction de I, p. 72-\uuline{75}
on voit que $f:X\mathrel{-}Y$ est dans $W$ ssi $f^\circ:X^\circ\mathrel{-}Y^\circ$
l'est. Donc L1,L3,L4 impliquent L3$'$, L4$'$
(axiomes duels de L3, L4).

Il y avait des axiomes complémentaires

L a) $f\in W\Longrightarrow\pi_0(f)$ bijectif.

L b) Le foncteur canonique de localisation
\[
\mathrm{Cat}\mathrel{\unsure{-|\to}}(\mathrm{Cat})W^{-1}\overset{\mathrm{déf}}{=}\mathrm{Hot}_W
\]
commute aux sommes finies

L c) Item pour produits finis.

Il faudrait quand-même dégager des conditions
concrètes sur $W$ pour la validité de L b, L c.
Les conditions L 3 quater et L \unsure{a}) s'écrivent
respectivement
\[
W_\infty\subset W,\qquad W\subset W_0
\]
\scanpage{59}
\mspage{3}
% archive sheet number: 58
Leur conjonction est l'axiome des plus raisonnables
\[
W_\infty\subset W\subset W_0\ .
\]
Bien sûr, l'axiome L 5 a maintenant l'air un
peu étriqué, i.e. faiblard. On va voir très vite de
quoi on aura besoin ! \deleted{\unsure{Si}}

On veut donc définir un dérivateur
\[
\mathrm{HOT}_W
\]
à partir de la catégorie de modèles
\[
(\mathrm{Cat},W\subset\mathrm{Fl}(\mathrm{Cat})).
\]

Donc pour définir
\[
\mathrm{HOT}_W(S)\qquad (S\in\mathrm{Ob}\,\mathrm{Cat}),
\]
la recette générale est de prendre la catégorie
\[
M(S)=\mathrm{Cat}(S)=\underline{\mathrm{Hom}}(S^\circ,(\mathrm{Cat})),
\]
d'y introduire l'ens de flèches
\[
\left\{\begin{aligned}
&W_S\subset\mathrm{Fl}(M(S))\\
&u\in W_S\ \Longleftrightarrow\ u_s\in W\quad\forall s\in\mathrm{Ob}\,S\ .
\end{aligned}\right.
\]
On \unsure{doit|choisit} poser
\[
\mathrm{HOT}_W(S)\overset{\mathrm{def}}{=}\underline{\mathrm{Hom}}(S^\circ,\mathrm{Cat})W_S^{-1}.
\]
Mais la catégorie $\underline{\mathrm{Hom}}(S^\circ,\mathrm{Cat})$ s'interprète,
à \deleted{équivalence} \inserted{isomorphisme \unsure{can.|de c.|de cat.}} près, comme celle
\[
\underline{\mathrm{Fib.}}\underline{\mathrm{sc.}}(S)
\]
des catégories $X$ \deleted{fib.} fibrées sur $S$, munies
d'un ``scindage'' i.e. d'un choix des foncteurs
changement de base $u^*:X_t\longrightarrow X_s$ (pour $t\xrightarrow{u}s$
dans $S$), avec de plus \uline{transitivité} \uline{stricte}
\[
(vu)^*=u^*v^*\qquad\text{si}\quad t\mathrel{\overset{u}{-}}s\xrightarrow{v}r\text{ dans }S,
\]
\scanpage{60}
\mspage{4}
% archive sheet number: 59
et non seulement transitivité à isom. can. près.
Ainsi on peut écrire aussi
\[
\mathrm{HOT}_W(S)\xrightarrow{\sim}\underline{\mathrm{Fib.sc.}}(S)W_S^{-1},\tag*{($\unsure{*|+}$)}
\]
où cette fois $W_S$ est interprété comme
\[
W_S\subset\mathrm{Fl}\bigl(\underline{\mathrm{Fib}}\,\underline{\mathrm{sc}}.(S)\bigr),
\]
savoir
\[
W_S=\{u\in\mathrm{Fl}(-)\mid u_s\in W\quad\substack{\text{pour tt}\\s\in\mathrm{Ob}\,S}\}.
\]
(foncteurs strictement cartésiens, i.e. commutant
aux \unsure{scindages|scindage}, qui induisent des $W$-équivalences
entre les catégories fibres).

La flèche de \unsure{(*)|cas} est bel et bien un iso.
de catégories, pas seulement une équivalence.
Mais on va définir d'autres catégories, plus
commodes à certains égards, qui seront
\uline{équivalentes} (on l'espère), mais sûrement
pas isomorphes, à $\mathrm{HOT}_W(S)$. (On s'occupera
plus tard de la variance de $\mathrm{HOT}_W(S)$ en $S$.)

Les catégories $\underline{\mathrm{Hom}}(S^\circ,\mathrm{Cat})$, $\underline{\mathrm{Fib.sc.}}(S)$,
ainsi que les variantes que nous allons
introduire maintenant, sont en fait des
2-catégories. Mais pour décrire $\mathrm{HOT}_W(S)$,
seules les structures de 1-Catégories sont utilisées. On a l'impression que c'est bien
\scanpage{61}
\mspage{5}
% archive sheet number: 60
\unsure{grossier}, qu'on perd quelque chose ce faisant.
Mais ce qui est perdu va se retrouver, et
surabondamment, dans la structure des
$\underline{\mathrm{Hom}}$ internes dans les $\mathrm{HOT}(S)$, si on se
``rappelle'' que ces $\underline{\mathrm{Hom}}$ internes, donnant
lieu à des ``types d'homotopie relatifs''
$\underline{\mathrm{Hom}}_{\mathrm{HOT}_W}(S;\unsure{\xi|\zeta},y)$ (si $\xi,y$ en sont), peuvent
s'interpréter en termes de $\infty$-champs \unsure{de|en} \deleted{\unsure{$S$}}
groupoïdes sur $S$, avec toute la richesse de
structure qui leur appartient. Mais ceci est
une autre histoire \ldots

En fait, je vois un diagramme de
\deleted{\unsure{9}}13 catégories de ``catégories sur $S$''\unsure{,} à examiner
à présent
\diagram{\begin{tikzcd}[ampersand replacement=\&,column sep=1pt,row sep=2ex,cells={font=\scriptsize}]
\deleted{\unsure{\mathrm{Cat}\,\mathrm{fib}}}
\&\underline{\mathrm{Fibsc}}_W(S)\arrow[dl,hook]\arrow[d,"\mathrm{fid}"]
\&\&\mathrm{Cofibsc}_W(S)\arrow[dr,hook]\arrow[d]
\&\\
\underline{\mathrm{Fibsc}}(S)\arrow[d,"\mathrm{fid}"]
\&\underline{\mathrm{Fib}}_W(S)\arrow[dl,hook]
\&\&\underline{\mathrm{Cofib}}_W(S)\arrow[dr,hook]
\&\underline{\mathrm{Cofibsc}}(S)\arrow[d,"\mathrm{fid}."]\\
\underline{\mathrm{Fib}}(S)\arrow[dr,hook]
\&\deleted{\unsure{\underline{\mathrm{Liss}}(S)\to}}
\&\underline{\unsure{\mathrm{Parf}|\mathrm{Perf}}}(S)\arrow[dl,hook]\arrow[d,hook]\arrow[dr,hook]
\&\deleted{\unsure{\underline{\mathrm{Prop}}(S)\to}}
\&\underline{\mathrm{Cofib}}(S)\arrow[dl,hook]\\
\&\underline{\mathrm{Liss}}\,S\arrow[dr,hook]
\&\underline{\mathrm{Fibhot}}(S)\arrow[d,hook]
\&\underline{\mathrm{Prop}}(S)\arrow[dl,hook]\&\\
\&\&\mathrm{Cat}_{/S},\&\&
\end{tikzcd}}
où $\underline{\mathrm{Fib}}$, $\underline{\mathrm{Cofib}}$ désignent les catégories Cat\unsure{-|. }fibrées,
Cat-cofibrées sur $S$, $\underline{\mathrm{Liss}}$ \deleted{\unsure{$S$}} et $\underline{\mathrm{Prop}}$ les catégories
\scanpage{62}
\mspage{6}
% archive sheet number: 61
lisses resp. propres sur $S$, $\underline{\unsure{\mathrm{Parf}|\mathrm{Perf}}}$ celles qui sont à la
fois lisses et propres (fibrations parfaites),
enfin $\underline{\mathrm{Fibhot}}$ les catégories Hot-fibrées sur $S$.
Toutes les flèches sont des foncteurs \unsure{str.} pleinement fidèles, à l'exception des deux \inserted{(plutôt quatre)} flèches
verticales supérieures, qui sont fidèles sans plus.
Je viens d'en rajouter encore 4, de \unsure{catégorie|catégories} :
$\underline{\mathrm{Fib}\deleted{\unsure{o}}\mathrm{sc}}_0,\underline{\mathrm{Fib}}_0,\underline{\mathrm{Cofibsc}}_0,\underline{\mathrm{Cofib}}_0$ sont les sous-catégories
pleines des catégories correspondantes sans indice,
formées des catégories fibrées resp. cofibrées \inserted{$X$} sur
$S$, pour lesquelles les foncteurs de chgt. de base
resp. \unsure{cochgt|co-chgt} de base
\[
u^*:X_t\mathrel{\unsure{-|\to}}X_s,\qquad\text{resp.}\quad u_*:X_s\mathrel{-} X_t,
\]
pour $s\xrightarrow{u}t$ dans $S$, sont des $\mathrm{Hot}$-équivalences.
En termes de foncteurs
\[
S^\circ\xrightarrow{F}\mathrm{Cat}\qquad\text{resp.}\quad S^{\deleted{\unsure{\circ}}}\xrightarrow{F'}\mathrm{Cat},
\]
cela signifie que ce foncteur transforme
flèches quelconques de $S$ en éléments de
$W$ dans $\mathrm{Cat}$, \deleted{\unsure{donc}} i.e qu'on a un diagramme
commutatif
\diagram{\begin{tikzcd}[ampersand replacement=\&,row sep=large]
\deleted{\unsure{\Pi}}\,S^\circ\arrow[r,"F"]\arrow[d,"\mathrm{can}"']\&
\mathrm{Cat}\arrow[d,"\mathrm{can}"]\\
\begin{aligned}&\unsure{(\Pi|\Pi}\,S^\circ=S^\circ(\mathrm{Fl}\,S^\circ)^{-1}\\&\simeq(\Pi S)^\circ\end{aligned}\arrow[r,"\overline F"']\&
\mathrm{Hot}_W=(\mathrm{Cat})W^{-1}
\end{tikzcd}}
(et de $=$ pour $F'$), i.e qu'on a un ``système
\scanpage{63}
\mspage{7}
% archive sheet number: 62
local'' $\overline F$ sur $S$, à valeurs dans $\mathrm{Hot}_W$, défini à
l'aide de $F$. Je rappelle qu'on trouve également
un tel syst. local. dans le cas d'une catégorie
Hot-fibrée sur $S$ (et $=$ un système local
à valeurs dans $\mathrm{Hot}$ \unsure{\}} lui-même, au lieu de
$\mathrm{Hot}_W$ \unsure{\}} -- on utilise ici tacitement que
$W_\infty\subset W$, i.e. \fbox{L3\,$|$\,quater}$^*$).
% slanted
\margin{On va le supposer vrai dans \unsure{toute|\deleted{tout}} la suite}
Et il en est de $=$ a fortiori pour les objets de
$\underline{\unsure{\mathrm{Parf}|\mathrm{Perf}}}(S)$. Pour $\mathrm{Liss}\,S$ resp. $\mathrm{Prop}.\,S$, on
trouve simplement un foncteur
\[
S^\circ\xrightarrow{\overline F}\mathrm{Hot}\qquad\text{resp.}\quad S\mathrel{\overset{\overline{F'}}{\unsure{-|\to}}}\mathrm{Hot},
\]
d'où en composant avec $\mathrm{Hot}\mathrel{-}\mathrm{Hot}_W$, des
foncteurs à valeurs dans $\mathrm{Hot}_W$.

Dans chacune de ces 14 catégories, on
introduit un ensemble de flèches, qui
serviront à localiser, et on le désignera tjrs
par la même notation $W_S$. Dans \deleted{toutes}
ces catégories sauf seulement $\mathrm{Cat}_{/S}$, \deleted{\unsure{la}}
$W_S$ désigne la $W$-équivalence ``fibre par fibre''.
Dans $\mathrm{Cat}_{/S}$ d'autre part, ce sera la
$W$-équivalence sans plus. toutes les flèches
du diagramme \unsure{p.|ci} sont compatibles avec les
ensembles localisants de flèches.
\scanpage{64}
\mspage{8}
% archive sheet number: 63
\diagram{\begin{tikzcd}[ampersand replacement=\&,column sep=2pt,row sep=2.6ex,cells={font=\scriptsize}]
\&\underline{\mathrm{Fibsc}}_0\,S\arrow[dl,hook,"\alpha'"']\arrow[d]
\&\&\underline{\mathrm{Cofibsc}}_0\,S\arrow[d]\arrow[dr,hook,"\beta'"]\&\\
\underline{\mathrm{Fibsc}}\,S\arrow[d]\&
\underline{\mathrm{Fib}}_0\,S\arrow[dl,hook,"\alpha"']\arrow[ddr,dashed,"\underline{\unsure{\mathrm{Lis}|\mathrm{Liss}}}_0"']
\&{}\&\mathrm{Cofib}_0\,S\arrow[dr,hook,"\beta"]\arrow[ddl,dashed,"\underline{\mathrm{Prop}}_0"]
\&\underline{\mathrm{Cofibsc}}\,S\arrow[d]\\
\underline{\mathrm{Fib}}\,S\arrow[dr,hook]\&
\&\underline{\unsure{\mathrm{Parf}|\mathrm{Perf}}}\,S\arrow[dl,hook]\arrow[d,hook]\arrow[dr,hook]
\arrow[u,phantom,"\deleted{\unsure{\Phi}}"]
\&\&\mathrm{Cofib}\,S\arrow[dl,hook]\\
\&\underline{\mathrm{Liss}\,S}\arrow[dr,hook]\&
\underline{\mathrm{Fibhot}}\,S\arrow[d,hook]\arrow[uuul,dashed,bend left=6,"\unsure{\Phi_0}"']\arrow[uuur,dashed,bend right=6,"\unsure{\Psi_0}"]\&
\underline{\mathrm{Prop}}\,S\arrow[dl,hook]\&\\
\&\&\mathrm{Cat}_{/S}\arrow[uuull,dashed,bend left=50,"\Phi"]\arrow[uuurr,dashed,bend right=50,"\Psi"']\&\&
\end{tikzcd}}

Je suis quasiment sûr du

\uline{Théorème}. Tous les foncteurs précédents induisent
des équivalences entre les catégories localisées.
% slanted
\margin{C'est totalement \unsure{canulé|coulé|croulé}, cf mes commentaires pages 32 ff. Solution des perplexités pages 41 ff. ! Voir \unsure{surtout} le ``théorème-scholie'' récapitulatif p. 70-81.}

Comme il y a 16 foncteurs, ça a l'air
d'un exercice de patience. Le principe général
est qu'on arrive à définir des foncteurs en
sens inverse, \inserted{compatibles avec les $W_S$,} de telle façon que les composés
soient \deleted{\unsure{iso}} reliés aux foncteurs identiques par
des flèches qui sont dans $W_S$, de sorte
que \deleted{par} par passage aux localisés ces foncteurs
induisent des (équivalences) quasi-inverses des
foncteurs \unsure{déduits|déduit} de ceux du diagramme.
Tous ces foncteurs en sens opposé résultent,
essentiellement par composition, de trois foncteurs
bien familiers. Deux de ceux-ci sont
indiqués en pointillés sur le diagramme, comme $\Phi,\Psi$.
\scanpage{65}
\mspage{9}
% archive sheet number: 64
Le troisième est un foncteur canonique
\[
\unsure{e|c|C}:\mathrm{Cat}_{/S}\longrightarrow\underline{\unsure{\mathrm{Parf}|\mathrm{Perf}}}\,S,
\]
qu'on va expliciter également.

Le foncteur $\Phi$ (resp. $\Psi$) appliqué à $\underline{\mathrm{Fib}}_0\,S$
(resp. $\mathrm{Cofib}_0\,S$), \deleted{\unsure{donne}} \inserted{induit} un foncteur
\[
\underline{\mathrm{Fib}}_0\,S\longrightarrow\underline{\mathrm{Fibsc}}_0
\qquad\text{respt.}\quad
\underline{\mathrm{Cofib}}_0\,S\longrightarrow\underline{\mathrm{Cofibsc}}_0\,S
\tag*{$(*)$}
\]
(ce \uline{qu'il faudra vérifier}). Cela fournit toutes
les flèches ``inverses'' dont on a besoin, à l'exception seulement des $\alpha,\beta,\alpha',\beta'$. Mais si on a
prouvé, \unsure{par} la méthode indiquée, que toutes les
flèches du diagramme induisent des équivalences sur les catégories localisées, pour prouver
qu'il en est encore de $=$ pour $\deleted{\unsure{\alpha}},\beta,\alpha',\beta'$, il
suffit de le prouver pour les \inserted{deux} foncteurs :
\[
\underline{\mathrm{Fibsc}}_0\,\deleted{(}S\deleted{)}\longrightarrow\underline{\mathrm{Fib}}\,S,\qquad
\underline{\mathrm{Cofibsc}}_0\,S\longrightarrow\underline{\mathrm{Cofib}}\,S\,;
\]
\unsure{or|où} il reste à définir des foncteurs en sens
inverse
\[
\underline{\mathrm{Fib}}\,\deleted{(}S\deleted{)}\to\underline{\mathrm{Fibsc}}_0\,\deleted{(}S\deleted{)},\qquad
\underline{\mathrm{Cofib}}\,S\unsure{\text{---}|\to}\mathrm{Cofibsc}_0\,S,
\]
et comme les sources sont des sous-catégories
pleines de $\mathrm{Cat}_{/S}$, il suffira de définir
\[
\Phi_0:\mathrm{Cat}_{/S}\longrightarrow\underline{\mathrm{Fibsc}}_0\,S,\qquad
\Psi_0:\mathrm{Cat}_{/S}\longrightarrow\mathrm{Cofib}\deleted{\unsure{s}}\mathrm{sc}_0\,S.
\]
Pour ceci, on posera
\[
\Phi_0=\Phi\unsure{\circ|*}C,\qquad\Psi_0=\Psi\circ C,\tag*{$\unsure{(**)|(xx)}$}
\]
\scanpage{66}
\mspage{10}
% archive sheet number: 65
où $C:\mathrm{Cat}_{/S}\to\underline{\unsure{\mathrm{Parf}|\mathrm{Perf}}}\,S$,
\inserted{$\Phi$ et $\Psi$} sont les 3 foncteurs \inserted{``rétrogrades''} \unsure{\deleted{ment}ionnés|mentionnés} plus haut, et $\unsure{r|i}:\underline{\mathrm{Parf}}\,S\hookrightarrow\mathrm{Cat}_{/S}$ est l'inclusion.
Il faut donc vérifier que $\Phi_0,\Psi_0$ définis
par $(**)$, qui à priori \deleted{\unsure{prennent}} ont comme catégories but $\underline{\mathrm{Fibsc}}\,S$ et $\underline{\mathrm{Cofib}}\,\underline{\mathrm{sc}}\,S$, prennent \unsure{leurs|leur}
valeurs dans les sous-catégories avec indices $0$.

Je vais donner un énoncé d'\unsure{équivalence|équivalences}
unique, qui couvrira \inserted{(presque)} tous les \deleted{\unsure{8}} \deleted{\unsure{18}}16 cas à la
fois :

\uline{Corollaire} Considérons les diagrammes
\diagram{\begin{tikzcd}[ampersand replacement=\&,column sep=small,row sep=small]
\&\underline{\mathrm{Parf}}\,S\arrow[dd,"r"']\&\\
\underline{\mathrm{Fibsc}}_0\,S\arrow[dr,"p"']\&\&\underline{\mathrm{Cofibsc}}_0\,S\arrow[dl,dash,"q"]\\
\&\mathrm{Cat}_{/S}\&
\end{tikzcd}}
et soit $\deleted{\unsure{(}}\unsure{\mathcal F|F}\deleted{\unsure{,\Sigma)}}$ une \inserted{ss-}catégorie \inserted{pleine de} \deleted{\unsure{avec un} ``localisateur''
$\Sigma\subset\mathrm{Fl}(\mathcal F)$, et $f:\mathcal F\to\mathrm{Cat}_{/S}$ un foncteur \deleted{\unsure{pleinement fidèle}} compatible avec les localisateurs. Supposons qu'il
existe un $\mathrm{Cat}_{/S}$-foncteur $g$ soit de $\underline{\mathrm{Fibsc}}_0\,S$, soit
de $\underline{\mathrm{Cofibsc}}_0\,S$, \inserted{soit de $\underline{\mathrm{Parf}}\,S$,} dans $\mathcal F$, compatible avec les
localisateurs. Alors $f$ et $g$ induisent des
équivalences entre les catégories localisées.}
\inserted{$\mathrm{Cat}_{/S}$, telle que l'un des foncteurs $p,q,r$ se factorise par $\mathcal F$. Soit $W_{\mathcal F}$ induit sur $\mathcal F$ par $W_S$ de $\mathrm{Cat}_{/S}$. Alors
$\mathcal F\hookrightarrow\mathrm{Cat}_{/S}$ induit une équivalence sur les catégories}
\deleted{(Ce \deleted{\illegible{1}} \inserted{cor.} implique \inserted{déjà} que $p$ \deleted{\unsure{et $f$}} et $q$ induisent des équivalences,
comme on voit en prenant $\mathcal F=\mathrm{Cat}_{/S}$, $f=\text{foncteur identique}$.)}
\inserted{localisées, et il en est de $=$ de $p,q,r$.}

\deleted{\uline{\unsure{Dém.}} On va se borner au cas $\underline{\mathrm{Fibsc}}_0\,S$, et
il suffira de prouver que $f$ induit une équivalence.
Pour conclure qu'il en sera de même de $g$,
il suffit de prouver qu'il en est ainsi de $p$.
Mais ce sera le cas, en appliquant le cor. au cas
$\mathcal F=\underline{\mathrm{Fibsc}}_0\,S,\ \Sigma=W_S,\ f=p$,}
